\documentclass{beamer}
\mode<presentation>
\usepackage{amsmath}
\usepackage{amssymb}
%\usepackage{advdate}
\usepackage{adjustbox}
\usepackage{subcaption}
\usepackage{enumitem}
\usepackage{multicol}
\usepackage{mathtools}
\usepackage{listings}
\usepackage{url}
\def\UrlBreaks{\do\/\do-}
\usetheme{Boadilla}
\usecolortheme{lily}
\setbeamertemplate{footline}
{
  \leavevmode%
  \hbox{%
  \begin{beamercolorbox}[wd=\paperwidth,ht=2.25ex,dp=1ex,right]{author in head/foot}%
    \insertframenumber{} / \inserttotalframenumber\hspace*{2ex} 
  \end{beamercolorbox}}%
  \vskip0pt%
}
\setbeamertemplate{navigation symbols}{}

\providecommand{\nCr}[2]{\,^{#1}C_{#2}} % nCr
\providecommand{\nPr}[2]{\,^{#1}P_{#2}} % nPr
\providecommand{\mbf}{\mathbf}
\providecommand{\pr}[1]{\ensuremath{\Pr\left(#1\right)}}
\providecommand{\qfunc}[1]{\ensuremath{Q\left(#1\right)}}
\providecommand{\sbrak}[1]{\ensuremath{{}\left[#1\right]}}
\providecommand{\lsbrak}[1]{\ensuremath{{}\left[#1\right.}}
\providecommand{\rsbrak}[1]{\ensuremath{{}\left.#1\right]}}
\providecommand{\brak}[1]{\ensuremath{\left(#1\right)}}
\providecommand{\lbrak}[1]{\ensuremath{\left(#1\right.}}
\providecommand{\rbrak}[1]{\ensuremath{\left.#1\right)}}
\providecommand{\cbrak}[1]{\ensuremath{\left\{#1\right\}}}
\providecommand{\lcbrak}[1]{\ensuremath{\left\{#1\right.}}
\providecommand{\rcbrak}[1]{\ensuremath{\left.#1\right\}}}
\theoremstyle{remark}
\newtheorem{rem}{Remark}
\newcommand{\sgn}{\mathop{\mathrm{sgn}}}
\providecommand{\abs}[1]{\left\vert#1\right\vert}
\providecommand{\res}[1]{\Res\displaylimits_{#1}} 
\providecommand{\norm}[1]{\lVert#1\rVert}
\providecommand{\mtx}[1]{\mathbf{#1}}
\providecommand{\mean}[1]{E\left[ #1 \right]}
\providecommand{\fourier}{\overset{\mathcal{F}}{ \rightleftharpoons}}
%\providecommand{\hilbert}{\overset{\mathcal{H}}{ \rightleftharpoons}}
\providecommand{\system}{\overset{\mathcal{H}}{ \longleftrightarrow}}
	%\newcommand{\solution}[2]{\textbf{Solution:}{#1}}
%\newcommand{\solution}{\noindent \textbf{Solution: }}
\providecommand{\dec}[2]{\ensuremath{\overset{#1}{\underset{#2}{\gtrless}}}}
\newcommand{\myvec}[1]{\ensuremath{\begin{pmatrix}#1\end{pmatrix}}}
\newcommand{\mydet}[1]{\ensuremath{\begin{vmatrix}#1\end{vmatrix}}}
\let\vec\mathbf

\lstset{
%language=C,
frame=single, 
breaklines=true,
columns=fullflexible
}

\numberwithin{equation}{section}


\title{5.5.9}
\author{Rushil Shanmukha Srinivas \\EE25BTECH11057 \\ Electrical Enggineering ,\\IIT Hyderabad.}

\date{\today} 
\begin{document} 

\begin{frame}
\titlepage
\end{frame}

\section*{Outline}
\begin{frame}
\tableofcontents
\end{frame}
\section{Problem}
\begin{frame}
\frametitle{Problem Statement}
\textbf{Question} : 
Using elementary row transformations find the inverse of the matrix 
\begin{align}
\myvec{3&0&-1\\
       2&3&0\\
       0&4&1}
\end{align}
\end{frame}
\section{Solution}
\subsection{ Augmented Matrix and Inverse of a matrix}
\begin{frame}
\frametitle{Augmented Matrix and Inverse of a matrix}
%\framesubtitle{Literature}
\textbf{Solution:}

Given, \begin{align} 
\vec{A}=\myvec{3&0&-1\\
       2&3&0\\
       0&4&1}
       \end{align}
Forming the Augmented Matrix
\begin{align}
\myvec{3&0&-1 &\vrule& 1&0&0 \\
       2&3&0 &\vrule&  0&1&0 \\
       0&4&1 &\vrule&  0&0&1 }
\end{align}
Applying elementary row operations to find the inverse
\begin{align}
\myvec{3&0&-1 &\vrule& 1&0&0 \\
       2&3&0 &\vrule&  0&1&0 \\
       0&4&1 &\vrule&  0&0&1 } \xrightarrow{R_3 \longrightarrow 3R_3-4R_2}
       \myvec{3&0&-1 &\vrule& 1&0&0 \\
       2&3&0 &\vrule&  0&1&0 \\
       -8&0&3 &\vrule&  0&-4&3} 
       \end{align}
       \end{frame}
       \begin{frame}
       \begin{align}
       \xrightarrow{R_3 \longrightarrow 3R_3+8R_1}
       \myvec{3&0&-1 &\vrule& 1&0&0 \\
       2&3&0 &\vrule&  0&1&0 \\
       0&0&1 &\vrule&  8&-12&9} 
\end{align}
\begin{align}
 \myvec{3&0&-1 &\vrule& 1&0&0 \\
       2&3&0 &\vrule&  0&1&0 \\
       0&0&1 &\vrule&  8&-12&9} \xrightarrow{R_2 \longrightarrow 3R_2-2R_1}
    \myvec{3&0&-1 &\vrule& 1&0&0 \\
       0&9&2 &\vrule&  -2&3&0 \\
       0&0&1 &\vrule&  8&-12&9} 
       \end{align}
       \begin{align}
       \xrightarrow{R_2 \longrightarrow R_2-2R_3}
        \myvec{3&0&-1 &\vrule& 1&0&0 \\
       0&9&0 &\vrule&  -18&27&-18 \\
       0&0&1 &\vrule&  8&-12&9} 
\end{align}
\begin{align}
 \myvec{3&0&-1 &\vrule& 1&0&0 \\
       0&9&0 &\vrule&  -18&27&-18 \\
       0&0&1 &\vrule&  8&-12&9} \xrightarrow{R_2 \longrightarrow \frac{R_2}{9}}
        \myvec{3&0&-1 &\vrule& 1&0&0 \\
       0&1&0 &\vrule&  -2&3&-2 \\
       0&0&1 &\vrule&  8&-12&9} 
       \end{align}
       \end{frame}
       \begin{frame}
       \begin{align}
       \xrightarrow{R_1 \longrightarrow R_1+R_3}
        \myvec{3&0&0 &\vrule& 9&-12&9 \\
       0&1&0 &\vrule&  -2&3&-2 \\
       0&0&1 &\vrule&  8&-12&9} 
\end{align}
\begin{align}
\myvec{3&0&0 &\vrule& 9&-12&9 \\
       0&1&0 &\vrule&  -2&3&-2 \\
       0&0&1 &\vrule&  8&-12&9} \xrightarrow{R_1 \longrightarrow \frac{R_1}{3}}
       \myvec{1&0&0 &\vrule& 3&-4&3 \\
       0&1&0 &\vrule&  -2&3&-2 \\
       0&0&1 &\vrule&  8&-12&9} 
\end{align}
The right side part of the Augmented Matrix is $\vec{A}^{-1}$
\begin{align}
\vec{A}^{-1} = \myvec{3&-4&3\\
                      -2&3&-2\\
                      8&-12&9}
\end{align}

\end{frame}

\section{C Code}
\begin{frame}[fragile]
\frametitle{C Code }
\begin{lstlisting}[language=C]  
 
#include <stdio.h>
#include <math.h>
// Function to display an augmented matrix [A|I]
void display(float mat[SIZE][2*SIZE]) {
    for (int i = 0; i < SIZE; i++) {
        for (int j = 0; j < 2*SIZE; j++) {
            printf("%10.6f ", mat[i][j]);
        }
        printf("\n");
    }
    printf("\n");
}

int main() {
    // Augmented matrix [A|I] for A = {{3,0,-1},{2,3,0},{0,4,1}}
    \end{lstlisting}
    \end{frame}
    \begin{frame}[fragile]
    \begin{lstlisting}
    float mat[SIZE][2*SIZE] = {
        {3.0f, 0.0f, -1.0f, 1.0f, 0.0f, 0.0f},
        {2.0f, 3.0f,  0.0f, 0.0f, 1.0f, 0.0f},
        {0.0f, 4.0f,  1.0f, 0.0f, 0.0f, 1.0f}
    };
    printf("Initial augmented matrix [A | I]:\n");
    display(mat);
    // Gauss-Jordan elimination
    for (int i = 0; i < SIZE; i++) {
        // Pivot element
        float pivot = mat[i][i];
        if (fabs(pivot) < 1e-6) {
            printf("Matrix is singular and cannot be inverted.\n");
            return 1;
        }
        // Normalize pivot row
        for (int j = 0; j < 2*SIZE; j++) {
            mat[i][j] /= pivot;
        }
\end{lstlisting}
\end{frame}
\begin{frame}[fragile]
\begin{lstlisting}
        // Eliminate column entries in other rows
        for (int r = 0; r < SIZE; r++) {
            if (r != i) {
                float factor = mat[r][i];
                for (int c = 0; c < 2*SIZE; c++) {
                    mat[r][c] -= factor * mat[i][c];
                }}}
        printf("After step %d:\n", i+1);
        display(mat);
    }

    // Extract inverse from right half
    printf("Inverse matrix A^{-1}:\n");
    for (int i = 0; i < SIZE; i++) {
        for (int j = SIZE; j < 2*SIZE; j++) {
            // Round if close to integer
            float v = mat[i][j];
\end{lstlisting}
\end{frame}
\begin{frame}[fragile]
\begin{lstlisting}
            if (fabs(v - roundf(v)) < 1e-6f)
                printf("%6.0f ", roundf(v));
            else
                printf("%8.6f ", v);
        }
        printf("\n");
    }

    return 0;
}

\end{lstlisting}
\end{frame}

\section{Python Code}
\begin{frame}[fragile]
\frametitle{Python Code 1: call\_c.py}
\begin{lstlisting}
import ctypes as ct
import numpy as np
handc1 = ct.CDLL("./func.so")
handc1.inv.argtypes = [
    ct.POINTER(ct.c_double),
    ct.POINTER(ct.c_double),
    ct.POINTER(ct.c_double)
]
A = np.array([3 , 2 , 0 ], dtype =np.float64).reshape(-1,1)
B = np.array([0 , 3 , 4 ], dtype =np.float64).reshape(-1,1)
C = np.array([-1 , 0 , 1 ], dtype =np.float64).reshape(-1,1)
handc1.inv.restype = None
handc1.inv(
    A.ctypes.data_as(ct.POINTER(ct.c_double)),
    B.ctypes.data_as(ct.POINTER(ct.c_double)),
    C.ctypes.data_as(ct.POINTER(ct.c_double)))
\end{lstlisting}
\end{frame}

\begin{frame}[fragile]
\frametitle{Python Code 2}
\begin{lstlisting}[language=Python]

import numpy as np
import numpy.linalg as LA
A = np.array([[3,0,-1],
              [2,3,0],
              [0,4,1]])

A_inv = LA.inv(A)

print("Matrix A:")
print(A)
print("\nInverse of A:")
print(A_inv)


\end{lstlisting}
\end{frame}

\end{document}
  